\begin{frame}{세 원리 비교}
  \small
  \begin{tabular}{@{}p{0.22\textwidth}p{0.26\textwidth}p{0.26\textwidth}p{0.26\textwidth}@{}}
    \toprule
     & VAE (우도 기반) & GAN (적대적) & Diffusion (스코어 기반) \\
    \midrule
    잠재변수 & 연속 \textbf{벡터} & \textbf{없음}(노이즈를 직접 변환)
      & 연속(노이즈 \textbf{단계}) \\
    학습 방법 & ELBO \textbf{최대화}(역전파) & \textbf{min-max 게임}의
      균형 & 각 단계 \textbf{노이즈 예측} 오차 최소화 \\
    학습 안정성 & \textbf{비교적 안정적} & \textbf{불안정하기 쉬움}
      & \textbf{안정적} \\
    생성 속도 & \textbf{빠름}(한 번에) & \textbf{빠름}(한 번에)
      & \textbf{느림}(반복 필요) \\
    \bottomrule
  \end{tabular}
  \vskip0.8em
  {\footnotesize 참고: 4.4절 GMM도 같은 질문의 \textbf{이산 잠재변수}
  버전 — $z$가 클러스터 번호라 E-step이 닫힌 형태로 \textbf{정확히}
  계산되므로, ``정밀하지만 이산''으로 여기 추가하지 않음.}
\end{frame}
